\subsection{Steenrod local finiteness at good primes}

Locally finite unstable modules are the first level of the usual
Krull filtration; see \cite[pp.~519--520]{schwartz2001krull}.
The argument below is given in full and does not apply the additional
hypotheses of that source's Theorem~0.1 to monoids.  The literature
search and exact attribution checks are recorded separately; no
priority claim is made.

All monoids in this subsection are finite pointed monoids: their
absorbing zero is the basepoint.  Actions and smash products are
pointed.  A finite group $H$ is treated as $H_+$.  Derived balanced
smash products are realized two-sided bars in pointed spaces, or
their suspension spectra when a stable comparison is specified.

\begin{definition}\label{E:controlled-action}
Fix $p$.  A pointed CW space with commuting cellular actions of finitely many
finite monoids is \emph{controlled} if every joint actor image is a
CW subcomplex and has finite total $\mathbb F_p$ homology.  The identity
actors are included, so the space itself has finite total homology.
We use the usual mapping-cylinder and simplicial-replacement CW
models; the explicit constructions below preserve this condition.
Only finite total homology, not a finite number of cells, is required.
\end{definition}

\begin{lemma}\label{E:LF-unstable-extension}
Locally finite unstable modules over the mod-$p$ Steenrod algebra
are closed under submodules, quotients, and extensions.  Consequently
a finite sequence of pointed homotopy pushouts and cofibers of spaces
with locally finite cohomology again has locally finite cohomology.
\end{lemma}

\begin{proof}
Submodules and quotients are immediate.  In an extension
$0\to K\to E\to Q\to0$, it suffices to consider a homogeneous
$x\in E$.  If $x$ lies in $K$ there is nothing to prove.  Otherwise take a finite homogeneous basis
of the Steenrod orbit of its image in $Q$.  Choose homogeneous lifts,
including $x$, of this basis.  Instability permits only finitely many
elementary squares, or powers and the Bockstein, on these finitely
many lifts.  Subtract the corresponding linear combinations of lifts
in their images.  The resulting finitely many errors lie in $K$.
Their orbits span a finite-dimensional submodule.  Its sum with the
span of the chosen lifts is finite-dimensional and closed under all
elementary operations, hence under the Steenrod algebra.  This
contains the orbit of $x$.

The cohomology exact sequence of a pointed cofiber or homotopy
pushout expresses the middle cohomology as an extension of
subquotients of the other cohomologies and their suspensions.
Suspension commutes with stable operations and preserves local
finiteness.  The middle cohomology is an unstable module, so the
extension assertion just proved applies.  No extension-closure
assertion for arbitrary stable modules is being used.
\end{proof}

\begin{lemma}\label{E:principal-image-replacement}
Let $R$ act on a controlled pointed space $X$ on the left, commuting
with an external finite-monoid action.  For a two-sided ideal
$P\subseteq R$, let $\mathcal P(P)$ be the finite poset of principal
right ideals $aR$ with $a\in P$, ordered by inclusion, and put
\[
 \mathcal F_P X=\mathop{\operatorname{hocolim}}_{aR\in\mathcal P(P)}
                         aX.
\]
The zero image is the one-point space; the homotopy colimit is
pointed.  Equal principal right ideals give literally equal images.
Then the following hold.
\begin{enumerate}
\item $\mathcal F_R X\to X$ is an $R$-equivariant underlying
      homotopy equivalence, also equivariant for the external action.
\item Ideal inclusions induce CW-subcomplex inclusions between the
      chosen models.  Left multiplication sends the diagram at $aR$
      to that at $baR$, strictly functorially.
\item Every joint actor image in $\mathcal F_P X$ has finite total
      mod-$p$ homology.  Its intersection with $\mathcal F_T X$, for
      a smaller ideal $T\subseteq P$, also has finite total homology.
      Consequently all joint images in
      $\mathcal F_P X/\mathcal F_T X$ have finite total homology.
\item If $P^h=0$, the spaces
      $F^sX=\mathcal F_{P^s}X$ ($P^0=R$) form a finite filtration,
      with strict multiplication maps
      $P^r\wedge F^sX\to F^{r+s}X$ for all $r,s\geq0$.
\end{enumerate}
There is a right-action version using principal left ideals.
\end{lemma}

\begin{proof}
If $aR=bR$, write $a=bu$ and $b=av$; then $aX=bX$.
If $aR\subseteq bR$, the corresponding map is the literal inclusion
$aX\subseteq bX$.  The object $1R$ is terminal in $\mathcal P(R)$,
proving the first assertion.  Multiplication $b$ sends $aR$ to
$baR$ and $aX$ to $baX$ and respects inclusions.  These maps compose
strictly.  External actors act objectwise and commute with them.

Here is a precise image description, including degeneracies.
The mapping-cylinder model is assembled from the finitely many
strict chains $a_0R\subsetneq\cdots\subsetneq a_nR$ and the pieces
$\Delta^n_+\wedge a_0X$.  A joint actor consisting of $b\in R$ and
an external actor $d$ sends such a chain to
$ba_0R\subseteq\cdots\subseteq ba_nR$.  Repeated vertices are
removed by the usual degeneracy map of simplices.  On the coordinate
space its image is all of
\[
 b d(a_0X)=d(ba_0X).
\]
If different chains have the same resulting chain, this coordinate
image depends only on its initial right ideal and on $d$, not on the
chosen lift.  Thus the actor image is the finite mapping-cylinder
realization over the simplicial image of the finite indexing nerve,
with these full coordinate-image spaces.  This simplicial image is
closed under faces.  A first face uses the inclusion of the initial
coordinate image into the next one; the other faces and the collapsed
degeneracies are precisely the ordinary diagram identifications.
Every coordinate image has finite total homology by the hypothesis
on joint actors.  A finite cellular realization therefore has finite
total homology.

Intersection with $\mathcal F_T X$ selects the finite subdiagram of
those image chains all of whose vertices belong to $\mathcal P(T)$.
The same reasoning applies.  These are CW-subcomplexes in the
mapping-cylinder model.  In the quotient the actor image is its
previous image divided by that intersection, so the cofiber exact
sequence proves the quotient assertion.  This also proves the claimed
CW-subcomplex condition for all the images constructed here.

For the final assertion, use the same construction also at $s=0$,
rather than replacing $F^0X$ by $X$.  The formula
\[
 b\cdot(a_0R\subsetneq\cdots\subsetneq a_nR,x)
   =(ba_0R\subseteq\cdots\subseteq ba_nR,bx)
\]
has its image in $F^{r+s}X$ when $b\in P^r$ and the original chain
lies in $P^s$.  Associativity and all filtration comparisons are
strict, since they are multiplication and inclusion of actual
ideals.  For $s\geq h$ the only image is the basepoint.  The right
version reverses multiplication.
\end{proof}

\begin{lemma}\label{E:regular-controlled-localization}
Let $J$ be a minimal nonzero regular $\mathcal J$-class ideal of a
finite pointed monoid $R$, and let $A=R/J$.  A controlled left
$R$-space $Y$, with a commuting external action, has a controlled
actual left $A$-space $Y'$ and an equivariant underlying equivalence
\[
 Y'\simeq A\mathbin{\wedge_R^{\mathbb L}}Y.
\]
The same holds on the right.  This comparison preserves the external
actions and can therefore be used inside longer balanced bars.
\end{lemma}

\begin{proof}
Choose $e\in J$ idempotent and let $H$ be its maximal group.
Rees coordinates give a literal pointed bimodule bijection
\[
 Re\wedge_H eR\ \xrightarrow{\ \cong\ }\ J.
\]
Both side sets are free on the relevant $H$-side, so this is also
their derived balanced product.  The evaluation map therefore gives
an actual mapping-cone model
\[
 Z=\operatorname{hocofib}(Re\wedge_H eY\longrightarrow Y)
       \simeq A\mathbin{\wedge_R^{\mathbb L}}Y.
\]
The bar preserves pointed homotopy cofibers; $eR\wedge_R^{\mathbb L}Y
\simeq eY$ follows by retracting the representable $R$-module.
The cone is an actual $R$-space.

Every joint actor image in $Z$ is the cone of the corresponding
image map.  Its source is a finite wedge of copies of a joint image
in $eY$: left multiplication on $Re$ has as its nonzero image a
union of free right $H$-orbits.  Its target is a joint image in $Y$.
Thus these cone images are CW-subcomplexes and have finite total
mod-$p$ homology.

For every nonzero $j\in J$, the image $jRe$ is one free right
$H$-orbit, and evaluation
\[
 jRe\wedge_H eY\longrightarrow jY
\]
is a homeomorphism.  Explicitly choose $u\in jRe$ and $v\in eR$
with $vu=e$ and $uv$ a left identity on the row containing $j$.
Then $u:eY\to jY$ and $v:jY\to eY$ are inverse homeomorphisms.
Consequently each pointed image $jZ$ is contractible: it is the cone
of this homeomorphism.  The zero image is also contractible.

Apply the principal-image construction to $Z$.  The subspace
$\mathcal F_JZ\subseteq\mathcal F_RZ$ is pointed contractible,
because every object of its pointed homotopy-colimit diagram is
contractible.  It is $R$-invariant, and every $j\in J$ sends all of
$\mathcal F_RZ$ into it literally.  Hence
\[
 Y'=\mathcal F_RZ/\mathcal F_JZ
\]
is an actual $A$-space.  The quotient map is an equivariant underlying
equivalence, because it collapses a contractible CW-subcomplex; and
$\mathcal F_RZ\to Z$ is an equivariant underlying equivalence.
The image description in Lemma~\ref{E:principal-image-replacement}
shows that $Y'$ is controlled, including all external actors.
This supplies strict quotient actions without assuming that arbitrary
coherent actions on finite spectra have controlled images.
\end{proof}

\begin{theorem}\label{E:general-controlled-multibar}
Let $R_1,\ldots,R_k$ be finite pointed monoids, and let
$X_0,\ldots,X_k$ be controlled pointed spaces with the commuting
left and right actions needed to form
\[
 Z=X_0\mathbin{\wedge_{R_1}^{\mathbb L}}X_1
       \mathbin{\wedge_{R_2}^{\mathbb L}}\cdots
       \mathbin{\wedge_{R_k}^{\mathbb L}}X_k.
\]
If $p$ divides none of the maximal-subgroup orders of these monoids,
then $\widetilde H^*(Z;\mathbb F_p)$ is locally finite over the
Steenrod algebra.
\end{theorem}

\begin{proof}
Induct on the finite multiset of the numbers of nonzero elements in
the monoids $R_i$.  Use the well-founded order in which one number
may be replaced by any finite list of strictly smaller numbers.
For example the ordinal sum $\sum_i\omega^{|R_i|-1}$, written in
decreasing order of exponents, strictly decreases under each such
replacement.  With no balancing monoids the assertion is the
finite-total-homology hypothesis on the one remaining space.
If a balancing monoid has $0=1$, unitality forces each adjacent
pointed module to be the one-point space, so its full bar is trivial.

Choose one balancing monoid $R$, and operate on its two adjacent
controlled bimodule spaces before performing the other bars.  Choose
a minimal nonzero $\mathcal J$-class ideal $J$ of $R$.

First suppose $J$ is regular.  Put $A=R/J$ and choose $e,H$ as in
Lemma~\ref{E:regular-controlled-localization}.  The ideal factorization
gives a pointed homotopy-cofiber sequence of balanced bars, with
terms at the chosen position
\[
 X e\mathbin{\wedge_H^{\mathbb L}}eY
   \longrightarrow X\mathbin{\wedge_R^{\mathbb L}}Y
   \longrightarrow X'\mathbin{\wedge_A^{\mathbb L}}Y'.
\]
Indeed $A\wedge_R^{\mathbb L}A\simeq A$: tensoring the cofiber
$J\to R\to A$ with $A$ kills $J$, since $eA=0$.
Apply the controlled strictification lemma to both adjacent modules
in the last term.  The remaining external actions commute with
these constructions.  The quotient monoid is smaller, and its local
groups are unchanged outside $J$, so the induction hypothesis applies
to the last term of the full multibar.

The first term can be absorbed into one new controlled bimodule.
For a finite group $H$ it is
\[
 EH_+\wedge_H(Xe\wedge eY),
\]
with the diagonal action using inversion on one side.  Since
$p\nmid|H|$, exact coinvariants show that this space has finite total
mod-$p$ homology.  Every external actor image is the same homotopy
orbit construction on the corresponding actual image subspaces;
it has finite total homology for the same reason and is a CW
subcomplex in the bar model.  Thus it is controlled.  This removes
the chosen balancing monoid entirely, and the induction hypothesis
applies to the remaining shorter chain.  The cofiber sequence and
Lemma~\ref{E:LF-unstable-extension} prove the assertion in the regular
case.

Now suppose $J$ is nonregular.  Then $J^2=0$, as follows directly
from finiteness.  The set $J^2$ is an $R$-ideal contained in the
minimal nonzero ideal $J$.  If it were nonzero, it would equal $J$,
and hence $J^n=J$ for every $n$.  Choose a nonzero product of more
than $|J|$ elements of $J$.  All its prefix products are nonzero,
so two prefix products agree, say $p_i=p_j=p_i u$ with $i<j$ and
$u\in J$.  A sufficiently large idempotent power $u^\omega$ then
satisfies $p_i u^\omega=p_i\ne0$.  It is a nonzero idempotent in
$J$, contrary to nonregularity.  Thus this is the square-zero ideal
case.  Filter the chosen monoid by
\[
 R^0=R,\qquad R^1=J,\qquad R^2=0.
\]
Replace the two adjacent endpoint spaces by their principal-image
models, using right ideals for the left action and left ideals for
the right action.  Their zeroth levels are equivariantly equivalent
to the original spaces.  Their first levels use the ideal $J$, and
their second levels are the basepoint.  All these inclusions are
cofibrations.  The associated-graded endpoint spaces are the actual
quotients $F^s/F^{s+1}$.  They are controlled for the external actions
and the quotient $A=R/J$, by
Lemma~\ref{E:principal-image-replacement}; the $J$ action is literally
zero on each quotient because it raises filtration.

Filter the absolute multibar by the total of these endpoint weights
and the weights of the letters of the chosen $R$-bar.  In degree $n$
of that bar, no weight exceeds $n+2$.  Hence the suspension spectrum
of its filtration part $F^LZ$ has connectivity tending to infinity
with $L$, even after realization in the other bar directions.  In
particular every fixed-degree cohomology class of $Z$ is pulled back
from an actual cofiber $Z/F^LZ$ for sufficiently large $L$.

It remains to prove that each such finite-weight cofiber has locally
finite cohomology.  Associated graded is strong monoidal for the
convolution filtration: the finite pushout-product cube identifies
the quotient in weight $w$ with the wedge over weight partitions of
the smash products of the corresponding quotients.  It also commutes
with realization.  These assertions are statements about actual
cofiber maps of filtered spaces, or equivalently their suspension
spectra, not about a homology spectral sequence preserving operations.

The graded algebra has weight-zero part $A$ and positive part $J$
in weight one, with $J^2=0$.  Its unit and augmentation split by
weight.  Relative bar normalization over $A$ therefore computes
weight $w$ by a finite simplicial realization: a nondegenerate term
has at most $w$ positive $J$ factors.  Every such term is a longer
chain of balanced bars in which the one occurrence of $R$ has been
replaced by finitely many occurrences of the strictly smaller monoid
$A$.  The intervening $J$ factors are finite pointed $A$-bisets;
the endpoint graded pieces are the controlled spaces just described;
all other factors are unchanged.  The induction hypothesis makes
the cohomology of every term locally finite.  Finite realization,
including its positive suspension shifts, preserves local finiteness
by Lemma~\ref{E:LF-unstable-extension}.  Thus every graded weight has
locally finite cohomology.  The finitely many weights below $L$ give
the same conclusion for $Z/F^LZ$.

Finally, an arbitrary class of $Z$ pulls back from one of these
finite-weight cofibers.  Pullback commutes with every Steenrod
operation, so its orbit is finite-dimensional.  This completes the
well-founded induction.
\end{proof}

\begin{corollary}\label{E:general-good-prime-LF}
Let $M$ be a finite monoid.  At every prime not dividing the order
of any maximal subgroup of $M$, the mod-$p$ cohomology of $BM$ is
locally finite over the Steenrod algebra.  The same statement holds
for finite categories and their local subgroups.
\end{corollary}

\begin{proof}
Adjoin a new absorbing basepoint to $M$, obtaining the pointed monoid
$M_+$.  Let both endpoints be $S^0$, with all elements of $M$ acting
as the identity and the basepoint acting as zero.  These are finite
controlled spaces, and their pointed two-sided bar is $(BM)_+$.
Apply Theorem~\ref{E:general-controlled-multibar}.  The categorical
assertion follows from the stable consolidation construction in
Proposition~\ref{E:category-consolidation}, which preserves local
groups and shifts the entire cohomology with its stable operations.
\end{proof}

\begin{corollary}\label{E:Sigma-CP-infty-excluded}
The space $\Sigma\mathbb{CP}^{\infty}$ is not the classifying-space
homotopy type of a finite monoid or a finite category.
\end{corollary}

\begin{proof}
There are only finitely many primes dividing the local subgroup
orders of a finite category.  Choose any other prime $p$.  If
$u\in H^2(\mathbb{CP}^{\infty};\mathbb F_p)$ is the usual generator,
then the suspended class $\sigma u$ has an infinite Steenrod orbit.
At $p=2$ the successive operations
$\mathrm{Sq}^{2},\mathrm{Sq}^{4},\mathrm{Sq}^{8},\ldots$ give
$\sigma u^2,\sigma u^4,\sigma u^8,\ldots$.  At an odd prime the
successive powers $\mathcal P^1,\mathcal P^p,\mathcal P^{p^2},\ldots$
give $\sigma u^p,\sigma u^{p^2},\sigma u^{p^3},\ldots$.
They are nonzero in distinct degrees, contradicting
Corollary~\ref{E:general-good-prime-LF}.
\end{proof}

\begin{remark}\label{E:general-audit}
The complete proof above was independently checked by the root and a
second proof agent, including the strict image construction, the
pointed-space localization triangle, and the fixed-weight multibar
induction.  The finite cellular control in
\path{code/a_general_good_prime_control.py} checks the case where the
ordinary union of two contractible regular-ideal images is a circle,
while their pointed principal-ideal homotopy colimit is contractible.
It verifies every action composition, chain-map equation, and strict
quotient action in that example.  This is a corroborating finite
control, not the source of the all-order theorem.  The proof makes
no claim of literature priority and has no census-completeness
dependency.
\end{remark}
